\section{Experimental Evaluation}

\subsection{Experimental Setup}

\subsubsection{Dataset and Model}

We conduct experiments on the CIFAR-10 dataset, which consists of 60,000 32×32 color images across 10 classes (50,000 training, 10,000 test). For Shapley value computation, we construct a global validation set by sampling from the validation split (10\%) of each client's local data, collecting up to 1,000 samples with a fixed random seed for reproducibility. This avoids using the test set for utility evaluation, preventing potential data leakage that could bias client selection toward test-set performance. We employ a convolutional neural network (CNN) with the following architecture:
\begin{itemize}
\item Two convolutional layers (16 and 32 filters, $3 \times 3$ kernels, padding 1) with batch normalization
\item Two max-pooling layers ($2 \times 2$)
\item Two fully connected layers ($2048 \to 128 \to 10$)
\item ReLU activation and dropout (0.5)
\end{itemize}

\subsubsection{Non-IID Data Partitioning}

To simulate realistic heterogeneous scenarios, we partition the training data across clients using the Dirichlet distribution with concentration parameter $\alpha = 0.1$, creating extreme non-IID conditions where each client's data is highly skewed toward a few classes. This setting is significantly more challenging than typical non-IID benchmarks ($\alpha \geq 0.5$).

\subsubsection{Federated Learning Parameters}

\begin{table}[h]
\centering
\caption{Federated Learning Configuration}
\begin{tabular}{ll}
\hline
\textbf{Parameter} & \textbf{Value} \\
\hline
Total clients ($N$) & 100 \\
Selected per round ($K$) & 10 \\
Training rounds ($T$) & 100 \\
Local epochs ($E$) & 2 \\
Batch size & 32 \\
Learning rate ($\eta$) & 0.01 \\
Optimizer & SGD with momentum (0.5) \\
\hline
\end{tabular}
\end{table}

\subsubsection{Energy and Channel Parameters}

\begin{table}[h]
\centering
\caption{Energy and Channel Configuration}
\begin{tabular}{ll}
\hline
\textbf{Parameter} & \textbf{Value} \\
\hline
Initial energy ($E_{\text{init}}$) & 500 \\
Energy threshold ($E_{\text{threshold}}$) & 50 \\
Energy budget per round ($\bar{E}$) & 2.0 \\
Noise power ($\sigma^2$) & 1.0 \\
Transmission energy cap ($E_{\max}$) & 50 \\
CPU capacitance ($\kappa$) & $10^{-28}$ \\
CPU frequency ($f$) & $10^{9}$ Hz \\
Cycles per sample ($C_{\text{cyc}}$) & 20 \\
Channel model & Rayleigh fading \\
\hline
\end{tabular}
\end{table}

\subsubsection{Dual Scheduling Parameters}

\begin{table}[h]
\centering
\caption{Dual Scheduling Configuration}
\begin{tabular}{ll}
\hline
\textbf{Parameter} & \textbf{Value} \\
\hline
Lyapunov control ($V$) & 10.0 \\
Shapley update method & Exponential smoothing ($\alpha=0.5$) \\
Shapley iterations ($M$) & 5 \\
Initial rounds (round-robin) & 10 \\
\hline
\end{tabular}
\end{table}

\subsubsection{Security Configuration}

\begin{table}[h]
\centering
\caption{AES-256-GCM Security Configuration}
\begin{tabular}{ll}
\hline
\textbf{Parameter} & \textbf{Value} \\
\hline
Encryption algorithm & AES-256-GCM \\
Key size & 256 bits \\
Nonce size & 96 bits (fresh per upload) \\
Authentication tag & 128 bits \\
Threat model & Honest-but-curious server \\
Plaintext retention & Ephemeral (destroyed post-use) \\
\hline
\end{tabular}
\end{table}

\subsubsection{Baseline Methods}

We compare our Dual Scheduling mechanism against the following baselines:

\begin{itemize}
\item \textbf{Random Selection}: Uniformly random client selection at each round (FedAvg baseline).
\item \textbf{Power of Choice (PoC)}: Two-stage selection that samples $d=30$ candidates and selects the top-$K=10$ clients with highest local loss values.
\item \textbf{Pure Shapley}: Client selection based solely on Shapley values without energy awareness.
\item \textbf{Pure Energy}: Client selection based solely on residual energy without contribution assessment.
\end{itemize}

All methods use the same CNN model, training hyperparameters, and non-IID data distribution for fair comparison.
